Krátké shrnutí role
AI může v kurzu sloužit jako personalizovaný tutor, který poskytuje adaptivní cvičení, cílené nápovědy a rychlou formativní zpětnou vazbu přizpůsobenou potřebám jednotlivého studenta.
Praktické ukázky komerčních nástrojů potvrzují, že rozhraní pro učitele umožňují sledovat, v jakých oblastech studenti chybují, a generovat pro ně upravené úlohy či doporučení (např. funkce „Pokrok v matematice“) \cite{kb_ai_101_tipu_pdf_04e4a115_page_8_1}.
Tyto funkce učitelům nabízí přehledy výkonu, trendy v čase a konkrétní doporučení pro intervenci.
Pro demonstrační a didaktické účely mohou posloužit i interaktivní světy (např. Minecraft Education) vhodné pro vizualizaci principů učení strojů a první praktické aktivity s žáky \cite{kb_ai_101_tipu_pdf_04e4a115_page_15_2}.
Interaktivní prostředí pomáhají studentům porozumět abstraktním konceptům prostřednictvím vizuálních a herních metafor.
Nástroje typu Microsoft 365 Copilot pak mohou usnadnit hromadnou generaci testů a kvízů, což učiteli uvolní kapacitu pro individuální podporu studentů \cite{kb_ai_101_tipu_pdf_04e4a115_page_36_1}.
Automatizace rutinních úkonů umožňuje věnovat více času reflexi, doplňujícím vysvětlením a práci s rizikovými studenty.

Režimy interakce (praktické varianty)
- Adaptivní cvičení a diagnostika: systém nabízí úlohy podle dosavadního výkonu studenta a učiteli zobrazuje přehled oblastí s chybami (viz „Pokrok v matematice“) \cite{kb_ai_101_tipu_pdf_04e4a115_page_8_1}.
- Adaptive režim může být nastaven tak, aby zvýšil opakování klíčových konceptů u studentů s nižší úspěšností.
- Hintování místo hotových řešení (general background knowledge): tutor dává postupné nápovědy, které vedou studenta k vlastnímu řešení místo poskytnutí kompletní odpovědi.
- Tento přístup podporuje rozvoj metakognitivních dovedností a samostatného myšlení.
- Sokratovský styl dotazování (general background knowledge): tutor klade vedoucí otázky zaměřené na vysvětlení myšlenkového postupu studenta.
- Otázky mohou být konstruovány tak, aby diagnostikovaly nesprávné předpoklady a pomohly je odhalit.
- Generování cvičných testů a materiálů pro opakování: Copilot/Forms lze použít pro rychlé vytvoření testů a jejich variant pro diferenciaci \cite{kb_ai_101_tipu_pdf_04e4a115_page_36_1}.
- Hromadné generování šablon zrychluje přípravu a umožňuje snadno vytvářet varianty pro různé úrovně.
- Demonstrace konceptů a motivace v herním prostředí: připravené světy v Minecraft Education lze využít pro praktické ukázky a motivaci studentů při seznámení s problematikou AI \cite{kb_ai_101_tipu_pdf_04e4a115_page_15_2}.
- Hry a simulace zvyšují angažovanost a dávají možnost bezpečného experimentování.

Praktická pravidla pro nastavení AI tutora (doporučené)
- Požadujte důkazy o postupu: při úkolech, kde je důležitý proces (např. matematika), zařaďte povinnost nahrát nebo odevzdat postup řešení; to zvyšuje ověřitelnost učení a umožní cílenou zpětnou vazbu \cite{kb_ai_101_tipu_pdf_04e4a115_page_8_1}.
- Požadavky na postup pomáhají odhalit pouhé kopírování odpovědí bez pochopení.
- Omezte generování kompletních řešení tam, kde je cílem rozvíjet postupy a dovednosti (general background knowledge).
- Místo toho nastavte systém, aby poskytoval nápovědu po krocích a kontrolní otázky.
- Transparentnost vůči studentům: jasně formulujte pravidla používání AI v kurzu (co je povoleno, jak deklarovat použití AI) — viz příslušné šablony v příručce (general background knowledge).
- Komunikujte také důsledky porušení pravidel a požadavky na citování použití nástrojů.
- Využívejte enterprise/školní varianty nástrojů tam, kde se zpracovávají osobní údaje či sledované artefakty studentů (general background knowledge); pro demonstrační aktivity je možné použít i volně dostupné učební světy (Minecraft Education) \cite{kb_ai_101_tipu_pdf_04e4a115_page_15_2}.
- Zvažte nastavení retence dat a přístupu ke záznamům, aby byla zajištěna ochrana soukromí.

Praktický workflow pro zavedení tutora na úrovni předmětu (krátký návod)
1. Definujte cíle tutorování: které dovednosti má tutor podporovat (např. procedurální řešení úloh, jazyková podpora, opakování pojmů). (general background knowledge)
2. Vyberte nástroje a režimy: pro procvičování matematiky zvažte nástroje s možností sběru postupu (viz „Pokrok v matematice“), pro tvorbu testů použijte Copilot/Forms pro hromadné generování otázek \cite{kb_ai_101_tipu_pdf_04e4a115_page_8_1,kb_ai_101_tipu_pdf_04e4a115_page_36_1}.
3. Nastavte pravidla interakce a omezení (hinty ano, kompletní řešení ne; požadavek na postup). Dokumentujte tato pravidla v sylabu (general background knowledge).
4. Pilotujte malou skupinu studentů: zaznamenejte artefakty práce, sbírejte zpětnou vazbu a upravte chování tutora podle zjištěných potřeb.
5. Pro ukázky principů použijte připravené hodiny v Minecraft Education tam, kde je vhodné vizuální či herní zpracování \cite{kb_ai_101_tipu_pdf_04e4a115_page_15_2}.
6. Škálování: rozšiřte adaptivní banku úloh a připravte šablony testů/kvízů pro různé úrovně pomocí Copilot/Forms \cite{kb_ai_101_tipu_pdf_04e4a115_page_36_1}.
- Během pilotu sledujte metriky zapojení, přesnosti a času stráveného u úloh a iterujte podle dat.

Příklad krátkých pravidel pro studenta (možno vložit do sylabu) — general background knowledge:
- Tutor může poskytovat nápovědy a kontrolovat postupy; kompletní řešení je zakázáno u zadaných úloh, kde se hodnotí postup.
- Při použití AI uvolněného typu (např. pro stylistickou pomoc) uveďte stručnou deklaraci použití a přiložte vlastní reflexi (max. 200 slov).
- Dodržujte zásady akademické poctivosti při využívání AI a konzultujte pochybné případy s vyučujícím.

Očekávané přínosy (konkrétní + obecné)
- Konkrétní: nástroje pro diagnostiku umožní učiteli snadno identifikovat oblasti, kde studenti chybují, a nasadit cílené intervence (viz „Pokrok v matematice“) \cite{kb_ai_101_tipu_pdf_04e4a115_page_8_1}.
- Obecné (general background knowledge): lepší personalizace, rychlejší formativní zpětná vazba a zvýšená angažovanost studentů při adaptivních aktivitách a interaktivních demonstracích (např. herní světy).
- Dlouhodobě může systém přispět ke zefektivnění výuky a lepšímu sledování pokroku meziročními daty.

Kontrolní seznam pro učitele (rychle)
- Má nastavené pravidlo pro sběr postupu u úloh? (viz \cite{kb_ai_101_tipu_pdf_04e4a115_page_8_1})
- Je vymezena politika pro povolené/zakázané chování tutora a je komunikována studentům? (general background knowledge)
- Je připraven demonstrační materiál (např. Minecraft Education world) pro první hodinu seznámení? (viz \cite{kb_ai_101_tipu_pdf_04e4a115_page_15_2})
- Jsou připraveny šablony pro rychlé generování testů a jejich varianty (Copilot/Forms)? (viz \cite{kb_ai_101_tipu_pdf_04e4a115_page_36_1})

Poznámka k bezpečnosti a ochraně dat
- Pokud tutor zpracovává osobní údaje nebo artefakty studentů, zvažte použití školní/enterprise varianty služby a postupy pro anonymizaci či omezení retence dat (general background knowledge).
- Vhodné demonstrace bez citlivých dat lze realizovat i pomocí bezplatných výukových světů a veřejných materiálů \cite{kb_ai_101_tipu_pdf_04e4a115_page_15_2}.
- Pravidelně kontrolujte nastavení sdílení a záloh, aby nedocházelo k nežádoucímu zpřístupnění citlivých informací.